\documentclass[11pt,letterpaper]{article}

% --- Packages ---
\usepackage[utf8]{inputenc}
\usepackage[T1]{fontenc}
\usepackage{amsmath,amssymb,amsthm}
\usepackage{algorithm}
\usepackage{algpseudocode}
\usepackage{geometry}
\usepackage[numbers,sort&compress]{natbib}
\usepackage{hyperref}
\usepackage{booktabs}
\usepackage{enumitem}
\usepackage{bm}

\geometry{margin=1in}
\hypersetup{colorlinks=true, linkcolor=blue, citecolor=blue, urlcolor=blue}

% --- Theorem environments ---
\newtheorem{definition}{Definition}[section]
\newtheorem{theorem}{Theorem}[section]
\newtheorem{proposition}{Proposition}[section]
\newtheorem{remark}{Remark}[section]

% --- Notation macros ---
\newcommand{\R}{\mathbb{R}}
\newcommand{\frm}{\mathbf{F}}
\newcommand{\pos}{\mathbf{p}}
\newcommand{\basis}{\mathbf{b}}
\newcommand{\std}{\mathbf{e}}
\newcommand{\eye}{\mathbf{I}}
\newcommand{\norm}[1]{\left\lVert #1 \right\rVert}
\newcommand{\code}[1]{\texttt{#1}}
\newcommand{\op}[1]{\textup{\textsc{#1}}}

\title{%
  \textbf{TurtleND: A Turtle That Carries Its Own Orthonormal Frame} \\[6pt]
  \large Frame-Based Navigation in Three and $N$ Dimensions \\
  by Givens Rotations
}

\author{
  Eric G. Suchanek, PhD \\
  \texttt{suchanek@flux-frontiers.com} \\[4pt]
  Flux-Frontiers \\
  \url{https://github.com/Flux-Frontiers/turtlend}
}

\date{\today}

\begin{document}
\maketitle

\begin{abstract}
A turtle is a moving point that carries its own coordinate system. It
moves along its heading and rotates its frame; it never refers to global
axes. This article describes \code{TurtleND}, a turtle that carries an
orthonormal frame of $N$ basis vectors through $\R^N$ and rotates by
Givens rotations in the plane of any two of them. The three-dimensional
turtle it generalizes, \code{Turtle3D}, was written in C in 1990 for
building protein structures from bond lengths, bond angles and dihedral
angles. The $N$-dimensional version was written for navigating learned
embedding spaces. We give the state, the operations and their
definitions, the invariants they preserve, the conventions that make
\code{TurtleND}$(3)$ agree with \code{Turtle3D} on every command, and
the algorithms for orienting the frame to arbitrary directions,
re-orthonormalizing it, and adding a dimension to a live turtle. The
implementation is a NumPy package, \code{turtlend}, of about five hundred
lines, with no other dependencies.
\end{abstract}

\tableofcontents
\newpage

%======================================================================
\section{Introduction}
\label{sec:intro}
%======================================================================

Turtle geometry, as Papert and later Abelson and diSessa presented it,
replaces the Cartesian description of a path with an intrinsic one
\citep{papert1980mindstorms, abelson1981turtle}. The turtle knows where
it is and which way it faces. It moves forward some distance and turns
through some angle, and a curve is the record of those commands. The
description is local: it does not depend on where the origin is or how
the axes are drawn.

The three-dimensional turtle adds a second and third direction to the
heading. With a heading, a left and an up vector, the turtle has a
complete right-handed frame, and the commands become \op{Move},
\op{Turn}, \op{Pitch} and \op{Roll}. This is exactly the
description a chemist uses for a molecule. A bond has a length, the
angle between two bonds is a bond angle, and the rotation about a bond
is a dihedral. A turtle placed on an atom can walk to the next atom by
one move, one yaw and one roll, with no reference to global coordinates
at all. I wrote such a turtle in C in 1990 for the Proteus modeling
program \citep{pabo1986computer}, ported it to Python for proteusPy
\citep{suchanek2024proteuspy, suchanek2026disulfide}, and it has built
thousands of disulfide bonds since.

The generalization to $N$ dimensions came from a different problem.
The WaveRider project \citep{suchanek2026waverider} navigates the
embedding spaces produced by neural networks, where $N$ is in the
hundreds or thousands, and where the data concentrate on a
low-dimensional manifold whose tangent space changes from place to
place. A turtle that carries $N$ orthonormal basis vectors, and that
can align them to the local principal directions of the data, gives a
moving coordinate system for such a space. The three-dimensional
rotations become instances of one primitive, a Givens rotation in the
plane of two basis vectors \citep{givens1958computation}.

Both turtles lived as duplicate copies inside proteusPy and WaveRider
until September 2026, when they were moved into a single package,
\code{turtlend} \citep{suchanek2026turtlendsoftware}. This article is
the description of that package. It covers the mathematics of the
frame and its rotations (Section~\ref{sec:prelim}), the
\code{TurtleND} state and operations (Section~\ref{sec:turtlend}), the
conventions that make it agree with \code{Turtle3D}
(Section~\ref{sec:compat}), the frame-construction algorithms
(Section~\ref{sec:algorithms}), the numerical behavior
(Section~\ref{sec:numerics}), and the two applications that motivated
it (Section~\ref{sec:applications}). It does not cover the manifold
navigation built on top of the turtle in WaveRider; that is the subject
of the WaveRider article.

%======================================================================
\section{Mathematical Preliminaries}
\label{sec:prelim}
%======================================================================

\begin{definition}[Orthonormal Frame]
\label{def:frame}
An \emph{orthonormal frame} in $\R^N$ is an ordered set of $N$ vectors
$\frm = (\basis_0, \basis_1, \ldots, \basis_{N-1})$ with
\begin{equation}
  \langle \basis_i, \basis_j \rangle = \delta_{ij}
  \qquad \forall\; i, j \in \{0, \ldots, N-1\},
\end{equation}
where $\delta_{ij}$ is the Kronecker delta. Equivalently, the matrix
$F \in \R^{N \times N}$ whose rows are the $\basis_i$ satisfies
$F F^\top = F^\top F = \eye_N$. We write $\std_i$ for the $i$-th
standard basis vector of $\R^N$, so that the identity frame is
$\frm = (\std_0, \ldots, \std_{N-1})$ and $F = \eye_N$.
\end{definition}

The row convention matters for what follows. With basis vectors as
rows, the product $F \mathbf{v}$ is the vector of inner products
$\langle \basis_i, \mathbf{v} \rangle$, which are the coordinates of
$\mathbf{v}$ in the frame, and $F^\top \bm{\ell}$ is the linear
combination $\sum_i \ell_i \basis_i$, which maps frame coordinates back
to $\R^N$.

\begin{definition}[Givens Rotation]
\label{def:givens}
A \emph{Givens rotation} $G(i, j, \theta)$ in $\R^N$ is the identity
outside the coordinate plane $(i, j)$ and a plane rotation inside it
\citep{givens1958computation, golub2013matrix}:
\begin{equation}
  G(i, j, \theta)
  = \eye_N
  + (\cos\theta - 1)\,(\std_i\std_i^\top + \std_j\std_j^\top)
  + \sin\theta\,(\std_i\std_j^\top - \std_j\std_i^\top).
  \label{eq:givens}
\end{equation}
Its only entries that differ from the identity are
$G_{ii} = G_{jj} = \cos\theta$, $G_{ij} = \sin\theta$ and
$G_{ji} = -\sin\theta$.
\end{definition}

A Givens rotation is orthogonal, $G G^\top = \eye_N$, and has
determinant $+1$. Givens introduced these rotations to reduce a matrix
to triangular form one plane at a time; the turtle uses them for the
opposite purpose, to rotate a frame one plane at a time while leaving
every other basis vector untouched.

\begin{proposition}[Frame Rotation]
\label{prop:rotation}
If $F$ is the row matrix of an orthonormal frame, then $F' = G(i, j,
\theta)\,F$ is the row matrix of an orthonormal frame, and its rows are
\begin{equation}
  \basis_i' = \cos\theta\,\basis_i + \sin\theta\,\basis_j, \qquad
  \basis_j' = \cos\theta\,\basis_j - \sin\theta\,\basis_i, \qquad
  \basis_k' = \basis_k \;\; (k \neq i, j).
  \label{eq:rowrot}
\end{equation}
\end{proposition}

\begin{proof}
$F' F'^\top = G F F^\top G^\top = G G^\top = \eye_N$. Rows $i$ and $j$ of
$GF$ are the stated combinations by Equation~\eqref{eq:givens}; every
other row of $G$ is a row of the identity.
\end{proof}

Equation~\eqref{eq:rowrot} is the whole of the rotation machinery. It
says that $\basis_i$ rotates toward $\basis_j$ by $\theta$ within the
plane the two span, $\basis_j$ follows it around, and nothing else
moves. In three dimensions, a rotation in the plane of two basis
vectors is a rotation about the third, which is why the named
three-dimensional rotations are special cases.

%======================================================================
\section{The TurtleND}
\label{sec:turtlend}
%======================================================================

\subsection{State and Conventions}

A \code{TurtleND} in $N \geq 2$ dimensions holds a position
$\pos \in \R^N$ and an orthonormal frame $\frm$ stored as the rows of
an $N \times N$ matrix. The basis vectors are named by index
(Table~\ref{tab:names}). A new turtle sits at the origin with the
identity frame: $\pos = \mathbf{0}$, $F = \eye_N$. The
\code{position}, \code{frame}, \code{heading}, \code{left}, \code{up}
and \code{basis(i)} accessors return copies, so the only way to change
the turtle is through its methods.

\begin{table}[h]
\centering
\caption{Frame conventions. The names are inherited from \code{Turtle3D}.}
\label{tab:names}
\begin{tabular}{@{}lll@{}}
\toprule
\textbf{Row} & \textbf{Name} & \textbf{Role} \\
\midrule
$\basis_0$ & heading & direction of \op{Move} \\
$\basis_1$ & left & \\
$\basis_2$ & up & requires $N \geq 3$ \\
$\basis_3, \ldots, \basis_{N-1}$ & unnamed & reached through \code{basis(i)} and \code{rotate} \\
\bottomrule
\end{tabular}
\end{table}

The turtle also carries a pen flag and an optional tape. When
\code{recording} is set, every \op{Move} appends the new position
to the tape, so a path can be recovered after the fact. The pen is kept
for compatibility with \code{Turtle3D} and has no effect on the
geometry.

\subsection{Operations}

Table~\ref{tab:ops} lists the operations and their definitions. All
angles in the public interface are in degrees; the implementation
converts to radians once, at the entry to \code{rotate}. A rotation
never changes the position and a move never changes the frame.

\begin{table}[h]
\centering
\caption{TurtleND operations. $G$ is the Givens rotation of
Definition~\ref{def:givens}, applied to the frame as in
Proposition~\ref{prop:rotation}.}
\label{tab:ops}
\begin{tabular}{@{}lll@{}}
\toprule
\textbf{Operation} & \textbf{Definition} & \textbf{Plane} \\
\midrule
$\op{Move}(d)$ & $\pos \leftarrow \pos + d\,\basis_0$ & \\
$\op{Rotate}(\theta, i, j)$ & $F \leftarrow G(i, j, \theta)\,F$ & $(\basis_i, \basis_j)$ \\
$\op{Turn}(\theta)$ & $\op{Rotate}(\theta, 0, 1)$ & heading, left \\
$\op{Roll}(\theta)$ & $\op{Rotate}(\theta, 1, 2)$ & left, up \\
$\op{Pitch}(\theta)$ & $\op{Rotate}(-\theta, 0, 2)$ & heading, up \\
$\op{Yaw}(\theta)$ & $\op{Rotate}(180^\circ - \theta, 0, 1)$ & heading, left \\
$\op{ToLocal}(\mathbf{g})$ & $\bm{\ell} = F\,(\mathbf{g} - \pos)$ & \\
$\op{ToGlobal}(\bm{\ell})$ & $\mathbf{g} = \pos + F^\top \bm{\ell}$ & \\
$\op{Orthonormalize}()$ & $F \leftarrow$ sign-preserving $Q$ of $F^\top = QR$ & \\
$\op{Orient}(\pos, \mathbf{h}, \mathbf{l})$ & Algorithm~\ref{alg:orient} & \\
$\op{OrientToward}(\mathbf{d})$ & Algorithm~\ref{alg:toward} & (heading, $\mathbf{d}$) \\
$\op{ExpandDim}()$ & Algorithm~\ref{alg:expand} & \\
\bottomrule
\end{tabular}
\end{table}

\op{Rotate} is the one primitive. \op{Turn}, \op{Roll},
\op{Pitch} and \op{Yaw} are calls to it with fixed plane
indices and, for \op{Pitch} and \op{Yaw}, a transformed angle.
The signs and the $180^\circ$ offset are not arbitrary; they are the
conventions of \code{Turtle3D}, and Section~\ref{sec:compat} explains
where they come from. \op{Roll} and \op{Pitch} require
$N \geq 3$ because they involve $\basis_2$. \op{Rotate} rejects
$i = j$ and indices outside $[0, N)$.

\subsection{Invariants}

\begin{proposition}[Frame Invariant]
\label{prop:frame}
\op{Rotate}, and therefore every named rotation, maps an
orthonormal frame to an orthonormal frame.
\end{proposition}

\begin{proof}
Proposition~\ref{prop:rotation}.
\end{proof}

\begin{proposition}[Coordinate Transform Invertibility]
\label{prop:coords}
For every $\mathbf{g} \in \R^N$ and every $\bm{\ell} \in \R^N$,
\begin{equation*}
\op{ToGlobal}(\op{ToLocal}(\mathbf{g})) = \mathbf{g}
\qquad\text{and}\qquad
\op{ToLocal}(\op{ToGlobal}(\bm{\ell})) = \bm{\ell}.
\end{equation*}
\end{proposition}

\begin{proof}
Let $\bm{\delta} = \mathbf{g} - \pos$. Then
$\op{ToGlobal}(F\bm{\delta}) = \pos + F^\top F \bm{\delta} = \pos +
\bm{\delta} = \mathbf{g}$, and
$\op{ToLocal}(\pos + F^\top\bm{\ell}) = F F^\top \bm{\ell} =
\bm{\ell}$, since $F^\top F = F F^\top = \eye_N$.
\end{proof}

\begin{proposition}[Move Commutes with the Frame]
\label{prop:move}
\op{Move}$(d)$ followed by \op{Rotate} gives the same position
as \op{Rotate} followed by \op{Move}$(d)$ if and only if the
rotation does not involve $\basis_0$.
\end{proposition}

\begin{proof}
\op{Move} adds $d\,\basis_0$ to the position. A rotation in a plane
not containing index $0$ leaves $\basis_0$ fixed, so the two orders add
the same vector. A rotation in a plane containing index $0$ changes
$\basis_0$, so the two orders add different vectors unless $d = 0$ or
$\theta$ is a multiple of $360^\circ$.
\end{proof}

Proposition~\ref{prop:move} is why the turtle's path is a sequence of
straight segments joined at the positions where a heading rotation
occurred. \op{Roll} never changes the path in any dimension;
\op{Turn}, \op{Pitch} and \op{Yaw} always do.

\subsection{The Local Coordinate System}

The frame gives the turtle a moving coordinate system centered on its
position. \op{ToLocal} expresses a global point in that system, one
inner product per basis vector, and \op{ToGlobal} is the inverse
(Proposition~\ref{prop:coords}). Both are single matrix-vector products
and cost $O(N^2)$.

Two uses recur. In molecular building, the turtle is placed on an atom
with its frame aligned to the local bond geometry, and the coordinates
of nearby atoms are read out in the turtle's frame. The same atom
positions expressed this way are the same for every residue of the same
type, regardless of where the residue sits in the protein; this is what
lets a residue template be built once and placed anywhere. In manifold
navigation, a gradient or a displacement is taken into the frame with
\op{ToLocal}, its components are modified independently (kept,
scaled or zeroed), and the result is taken back out with
\op{ToGlobal} \citep{suchanek2026waverider}. The frame is the
device that makes ``the component along the second principal direction''
a meaningful phrase.

%======================================================================
\section{Agreement with Turtle3D}
\label{sec:compat}
%======================================================================

\code{Turtle3D} stores three named vectors, heading $\mathbf{h}$, left
$\mathbf{l}$ and up $\mathbf{u}$, and implements each rotation as an
explicit pair of component updates. Its conventions are the ones that
were convenient for building molecules in 1990, and \code{TurtleND}
adopts them so that the two turtles are interchangeable in three
dimensions. Table~\ref{tab:conventions} lists the four rotations as
\code{Turtle3D} writes them and the Givens plane and angle that
reproduce each.

\begin{table}[h]
\centering
\caption{The \code{Turtle3D} rotations and their \code{TurtleND}
equivalents. $c = \cos\alpha$ and $s = \sin\alpha$ for the angle
$\alpha$ given in the third column.}
\label{tab:conventions}
\begin{tabular}{@{}llll@{}}
\toprule
\textbf{Turtle3D} & \textbf{Update} & \textbf{Angle} $\alpha$ & \textbf{TurtleND} \\
\midrule
\code{turn(a)}  & $\mathbf{h}' = c\mathbf{h} + s\mathbf{l}$, $\mathbf{l}' = c\mathbf{l} - s\mathbf{h}$ & $a$ & \code{rotate(a, 0, 1)} \\
\code{roll(a)}  & $\mathbf{l}' = c\mathbf{l} + s\mathbf{u}$, $\mathbf{u}' = c\mathbf{u} - s\mathbf{l}$ & $a$ & \code{rotate(a, 1, 2)} \\
\code{pitch(a)} & $\mathbf{h}' = c\mathbf{h} - s\mathbf{u}$, $\mathbf{u}' = c\mathbf{u} + s\mathbf{h}$ & $a$ & \code{rotate(-a, 0, 2)} \\
\code{yaw(a)}   & $\mathbf{h}' = c\mathbf{h} + s\mathbf{l}$, $\mathbf{l}' = c\mathbf{l} - s\mathbf{h}$ & $180^\circ - a$ & \code{rotate(180 - a, 0, 1)} \\
\bottomrule
\end{tabular}
\end{table}

Three of the four follow directly from Equation~\eqref{eq:rowrot}.
\code{turn} and \code{roll} are rotations of the first index toward the
second by the given angle. \code{pitch} rotates the heading toward
$-\mathbf{u}$, which is a rotation of $\basis_0$ toward $\basis_2$ by
$-a$. The fourth, \code{yaw}, is a turn by the supplement of the
angle, and it exists because of how bond angles are defined.

\begin{remark}[Why yaw is $180^\circ - \theta$]
\label{rem:yaw}
Place the turtle on atom $A$ heading toward atom $B$, and move the bond
length so that it stands on $B$. Its heading is now the direction
$A \to B$. The bond angle $\theta_{ABC}$ is the angle at $B$ between the
bonds $B \to A$ and $B \to C$, and $B \to A$ is the reverse of the
heading. To point the heading along $B \to C$, the turtle must turn by
$180^\circ - \theta_{ABC}$. \code{yaw} takes the bond angle from the
chemistry tables and does the subtraction, so that building an atom is
\code{move(length)}, \code{yaw(angle)}, \code{roll(dihedral)} with the
tabulated values and nothing else. Users who want a plain turn by
$\theta$ call \code{turn}.
\end{remark}

\begin{theorem}[Equivalence in Three Dimensions]
\label{thm:compat}
For every finite sequence of \op{Move}, \op{Turn},
\op{Roll}, \op{Pitch} and \op{Yaw} commands with the same
arguments, \code{TurtleND}$(3)$ started at the identity frame and
\code{Turtle3D} started at its default frame end at the same position
with the same heading, left and up vectors, up to floating-point
rounding.
\end{theorem}

\begin{proof}
Both turtles start at $\pos = \mathbf{0}$ with
$(\mathbf{h}, \mathbf{l}, \mathbf{u}) = (\std_0, \std_1, \std_2)$.
\op{Move} is the same update in both. For each rotation,
Table~\ref{tab:conventions} shows that the \code{Turtle3D} component
updates are Equation~\eqref{eq:rowrot} for the stated plane and angle,
which is what \code{TurtleND} computes. Both implementations renormalize
the two changed vectors after each rotation, and renormalization of
vectors that are already unit length up to rounding changes them only
by rounding. By induction on the length of the sequence the states
agree at every step.
\end{proof}

The test suite of \code{turtlend} checks this theorem over a random
sequence of 200 commands drawn from all five operations with arguments
uniform in $[-180, 180]$, comparing position, frame and both coordinate
transforms at the end of the sequence to a tolerance of $10^{-7}$, and
separately on every individual operation.

\begin{remark}[Handedness of \op{Orient}]
\label{rem:handed}
The theorem is stated for the five movement commands and not for
\op{Orient}. \code{Turtle3D.orient} completes the frame with a
cross product, $\mathbf{u} = \mathbf{h} \times \mathbf{l}$, and is
always right-handed. \code{TurtleND.orient} completes it by
Gram-Schmidt against the standard basis vectors
(Algorithm~\ref{alg:orient}), which has no notion of handedness in
general $N$, and in three dimensions can produce
$\basis_2 = -\mathbf{h} \times \mathbf{l}$. The frame is orthonormal
either way. A caller who needs the right-handed completion in three
dimensions should use \code{Turtle3D}, or check the sign of
$\langle \basis_2, \mathbf{h} \times \mathbf{l} \rangle$ and negate
$\basis_2$ if it is negative.
\end{remark}

%======================================================================
\section{Frame Construction Algorithms}
\label{sec:algorithms}
%======================================================================

The rotations move a frame that already exists. Three further
operations build or rebuild one.

\subsection{Orient}

\op{Orient} places the turtle at a point and aims its heading and
left at two other points. This is the operation that puts a turtle on
an atom: the position is the alpha carbon, the heading point is the
carbonyl carbon or the beta carbon, and the left point is the nitrogen.
The remaining $N - 2$ basis vectors are not determined by the inputs
and are completed by Gram-Schmidt from the standard basis
\citep{golub2013matrix}.

\begin{algorithm}[H]
\caption{\op{Orient}$(\pos, \mathbf{h}, \mathbf{l})$}
\label{alg:orient}
\begin{algorithmic}[1]
\Require Position $\pos$, a point $\mathbf{h}$ to head toward, a point $\mathbf{l}$ on the left
\Ensure Frame $F$ orthonormal with $\basis_0 \parallel \mathbf{h} - \pos$ and $\basis_1$ in the plane of $\basis_0$ and $\mathbf{l} - \pos$
\State $\basis_0 \leftarrow (\mathbf{h} - \pos) / \norm{\mathbf{h} - \pos}$
\State $\mathbf{v} \leftarrow (\mathbf{l} - \pos) - \langle \mathbf{l} - \pos, \basis_0 \rangle\,\basis_0$
\State $\basis_1 \leftarrow \mathbf{v} / \norm{\mathbf{v}}$
\For{$k = 2, \ldots, N-1$}
  \State $\mathbf{v} \leftarrow \std_k - \sum_{m < k} \langle \std_k, \basis_m \rangle\,\basis_m$
  \If{$\norm{\mathbf{v}} < 10^{-12}$} \Comment{$\std_k$ lies in the span of the earlier vectors}
    \State try $\std_0, \std_1, \ldots$ in turn until one gives $\norm{\mathbf{v}} \geq 10^{-12}$
  \EndIf
  \State $\basis_k \leftarrow \mathbf{v} / \norm{\mathbf{v}}$
\EndFor
\end{algorithmic}
\end{algorithm}

The fallback on line 7 is needed. If the heading is aimed along
$\std_2$, then at $k = 2$ the projection of $\std_2$ onto $\basis_0$
removes all of it and the residual is zero; another standard basis
vector must be used instead. The tests include this case. The
completion is deterministic, which is worth more than any particular
choice of the extra vectors: the same inputs always produce the same
frame.

\subsection{Orient Toward}

\op{OrientToward} rotates the heading to a given direction
$\mathbf{d}$ with a single Givens rotation, and reports the angle it
rotated through. In WaveRider this is called with a principal direction
from a local PCA, or with the negative time axis
(Section~\ref{sec:time}).

\begin{algorithm}[H]
\caption{\op{OrientToward}$(\mathbf{d})$}
\label{alg:toward}
\begin{algorithmic}[1]
\Require Direction $\mathbf{d} \in \R^N$
\Ensure $\basis_0 = \mathbf{d} / \norm{\mathbf{d}}$; returns the rotation angle in degrees
\If{$\norm{\mathbf{d}} < 10^{-12}$} \Return $0$ \EndIf
\State $\mathbf{d} \leftarrow \mathbf{d} / \norm{\mathbf{d}}$
\State $c \leftarrow \langle \basis_0, \mathbf{d} \rangle$, clipped to $[-1, 1]$
\If{$|c| > 1 - 10^{-10}$} \Return $0$ \EndIf \Comment{already parallel or antiparallel}
\State $\theta \leftarrow \arccos c$
\State $\mathbf{q} \leftarrow \mathbf{d} - c\,\basis_0$;\quad $\mathbf{q} \leftarrow \mathbf{q} / \norm{\mathbf{q}}$ \Comment{unit component of $\mathbf{d}$ perpendicular to heading}
\State $j \leftarrow \arg\max_{m \geq 1} |\langle \basis_m, \mathbf{q} \rangle|$ \Comment{basis vector most aligned with $\mathbf{q}$}
\State $\basis_j \leftarrow \mathbf{q}$ \Comment{temporary: frame is no longer orthonormal}
\State $F \leftarrow G(0, j, \theta)\,F$ \Comment{now $\basis_0 = c\,\basis_0 + \sin\theta\,\mathbf{q} = \mathbf{d}$}
\State \op{Orthonormalize}$()$
\State \Return $\theta$ in degrees
\end{algorithmic}
\end{algorithm}

The rotation is in the plane spanned by the current heading and
$\mathbf{d}$, which is the plane of $\basis_0$ and $\mathbf{q}$. That
plane is not in general spanned by any two basis vectors of the frame,
so the algorithm swaps $\mathbf{q}$ into the slot of the basis vector
it most resembles, rotates, and then repairs the frame. Because
$\basis_0$ comes first in the Gram-Schmidt order of
\op{Orthonormalize}, the repair leaves the new heading exactly
$\mathbf{d}$ and adjusts the other vectors. The antiparallel case
$c = -1$ returns without rotating; the plane of rotation is undefined
there, and the caller who wants the turtle turned around has a whole
$(N-1)$-sphere of choices and must make one.

\subsection{Orthonormalize}

Each rotation renormalizes the two basis vectors it changes but does
not restore their orthogonality to each other or to the rest of the
frame. \op{Orthonormalize} rebuilds the whole frame from a QR
decomposition of $F^\top$, whose $Q$ factor is the Gram-Schmidt
orthonormalization of the basis vectors in order $0, 1, \ldots, N-1$
\citep{golub2013matrix}. A QR routine is free to return $-\mathbf{q}_k$
in place of $\mathbf{q}_k$, so each column is compared with the basis
vector it replaces and negated if the inner product is negative. The
result is the orthonormal frame nearest to the drifted one in the sense
that $\basis_0$ keeps its direction exactly, $\basis_1$ keeps its
component orthogonal to $\basis_0$, and so on down the list.

\subsection{Expand Dimension}
\label{sec:expand}

\op{ExpandDim} grows a live turtle from $N$ to $N + 1$ dimensions
without disturbing anything it has done so far. The position gains a
zero coordinate, each existing basis vector gains a zero component, and
the new basis vector is the new standard axis $\std_N$.

\begin{algorithm}[H]
\caption{\op{ExpandDim}$()$}
\label{alg:expand}
\begin{algorithmic}[1]
\Ensure Turtle in $\R^{N+1}$; returns the index $N$ of the new axis
\State $\pos \leftarrow (\pos, 0)$
\State $F \leftarrow \begin{pmatrix} F & \mathbf{0} \\ \mathbf{0}^\top & 1 \end{pmatrix}$
\State $N \leftarrow N + 1$
\State \Return $N - 1$
\end{algorithmic}
\end{algorithm}

\begin{proposition}
\label{prop:expand}
\op{ExpandDim} maps an orthonormal frame to an orthonormal frame,
and the original $N$ basis vectors are unchanged in their first $N$
components.
\end{proposition}

\begin{proof}
The new matrix is block diagonal with orthogonal blocks $F$ and $(1)$,
so it is orthogonal. The upper-left block is $F$.
\end{proof}

The new axis is the only basis vector with a nonzero last component,
so every existing rotation and move of the turtle takes place in the
original $N$-dimensional subspace until a rotation involving the new
index is issued. The operation is $O(N^2)$ and is a copy, not a
recomputation.

%======================================================================
\section{Numerical Behavior}
\label{sec:numerics}
%======================================================================

\subsection{Cost}

Table~\ref{tab:cost} gives the cost of each operation. A rotation
touches two rows of the frame and costs $O(N)$; this is the property
that makes the Givens formulation attractive at large $N$, since a
general rotation matrix applied to the frame would cost $O(N^3)$. The
coordinate transforms are one matrix-vector product each.

\begin{table}[h]
\centering
\caption{Cost per operation for a turtle in $N$ dimensions.}
\label{tab:cost}
\begin{tabular}{@{}lll@{}}
\toprule
\textbf{Operation} & \textbf{Cost} & \textbf{Notes} \\
\midrule
\op{Move} & $O(N)$ & \\
\op{Rotate} and named rotations & $O(N)$ & two rows, two norms \\
\op{ToLocal}, \op{ToGlobal} & $O(N^2)$ & one matrix-vector product \\
\op{Orient} & $O(N^3)$ & Gram-Schmidt over $N - 2$ vectors \\
\op{OrientToward} & $O(N^3)$ & one rotation, one \op{Orthonormalize} \\
\op{Orthonormalize} & $O(N^3)$ & QR of $F^\top$ \\
\op{ExpandDim} & $O(N^2)$ & copy into a larger array \\
\bottomrule
\end{tabular}
\end{table}

\subsection{Drift}

In exact arithmetic the frame is orthonormal forever. In floating
point each rotation introduces rounding of order the machine epsilon,
about $2 \times 10^{-16}$ for double precision, into the two rows it
changes. Renormalizing those rows removes the error in their lengths
and leaves the error in their mutual angle, and in their angles with
the rest of the frame, to accumulate. The accumulation is a random
walk, so after $M$ rotations the departure from orthogonality is of
order $\sqrt{M}\,\epsilon$ in the typical case and $M\epsilon$ in the
worst case. At $M = 10^6$ rotations that is between $10^{-13}$ and
$10^{-10}$. \code{Turtle3D} has built molecules for thirty-six years
with this scheme and no orthonormalization step at all, because a
molecule takes a few hundred rotations. A long navigation in $N$
dimensions can call \op{Orthonormalize} every few thousand steps
at a cost that is negligible against a single PCA of the same size.

\subsection{Determinism}

Every operation is deterministic. \op{Orient} completes the frame
from the standard basis in a fixed order, \op{OrientToward} picks
its rotation plane by a deterministic argmax, and
\op{Orthonormalize} resolves the QR sign ambiguity against the
prior frame. No random numbers are drawn anywhere in the package. A
sequence of commands on a given start state produces the same frame on
every run, which is a requirement for the reproducibility of anything
built on top.

%======================================================================
\section{Applications}
\label{sec:applications}
%======================================================================

\subsection{Building Molecules}
\label{sec:molecules}

The original purpose of the turtle was to build protein structures
from internal coordinates. Each atom after the first three is placed by
three numbers: the length of the bond that attaches it, the angle that
bond makes with the previous one, and the dihedral angle about the
previous bond. A turtle standing on the previous atom, heading along
the previous bond, with its left vector in the plane of the previous
two bonds, places the new atom by
\begin{center}
\code{roll(dihedral)}, \quad \code{yaw(bond\_angle)}, \quad \code{move(bond\_length)},
\end{center}
and is then standing on the new atom, ready for the next
(Remark~\ref{rem:yaw}). No global coordinates enter at any point. The
same three commands build a peptide backbone, a side chain, or a
disulfide bridge, and the shape of the molecule is set by the dihedral
angles alone once the lengths and angles are fixed at their
tabulated values.

\code{Turtle3D} adds two conventions for proteins. A turtle on an alpha
carbon can be in the \emph{backbone} orientation, heading toward the
carbonyl carbon with the nitrogen on its left, or the \emph{side-chain}
orientation, heading toward the beta carbon with the nitrogen on its
left. \code{orient\_at\_residue} places the turtle in either
orientation from the atoms of a residue, and \code{bbone\_to\_schain}
and \code{schain\_to\_bbone} convert between them with a fixed sequence
of four rotations, since the tetrahedral geometry at the alpha carbon is
the same for every residue. proteusPy uses this to build a disulfide
bond in the turtle's local frame from its five dihedral angles and then
transform it into the frame of any residue in the Protein Data Bank
\citep{suchanek2024proteuspy, suchanek2026disulfide}.

\subsection{Navigating an Embedding Space}
\label{sec:navigation}

WaveRider uses \code{TurtleND} as the moving coordinate system for a
walk on a data manifold in $\R^N$ \citep{suchanek2026waverider}. At
each position it takes the $k$ nearest embedding vectors, computes
their covariance, and reads off the principal directions and
eigenvalues. The frame is aligned so that $\basis_0$ is the first
principal direction, $\basis_1$ the second, and so on; the number $d$
of eigenvalues needed to reach a fixed fraction of the variance is the
local intrinsic dimensionality. A step is then taken by bringing the
descent direction into the frame with \op{ToLocal}, zeroing the
components with index $d$ and above, scaling the rest by their
eigenvalues, and returning with \op{ToGlobal}. The turtle supplies
the frame, the transforms, and the rotations that steer between
principal directions; the covariance, the thresholds and the objective
belong to the walker, which is why they are not in this package. The
rest of that method, and its use in estimating intrinsic dimensionality
and sizing classifiers, is in the WaveRider article.

\subsection{Adding a Time Axis}
\label{sec:time}

\op{ExpandDim} exists for a corpus that has both a semantic
embedding and a date. Embedding the text gives a turtle in $\R^N$;
appending a standardized fractional year as coordinate $N$ gives a
turtle in $\R^{N+1}$ that can be steered in time as easily as in
meaning. \code{orient\_in\_time} is the convenience call
$\op{OrientToward}(-\std_N)$. The sign is a convention of the
corpus it was written for, whose time coordinate is a $z$-score that
decreases toward the present, so that ``forward in time'' is the
negative direction of the axis. Callers with the opposite convention
pass $+\std_N$ to \op{OrientToward} directly.

%======================================================================
\section{Implementation}
\label{sec:implementation}
%======================================================================

\code{turtlend} is a Python package on PyPI with a single runtime
dependency, NumPy \citep{harris2020numpy}. It provides three classes.
\code{TurtleND} is the $N$-dimensional turtle of
Section~\ref{sec:turtlend}; its frame is an $N \times N$ NumPy array of
doubles with basis vectors as rows, and \op{Rotate} is implemented
as Equation~\eqref{eq:rowrot} on two rows followed by renormalization.
\code{Turtle3D} is the three-dimensional turtle with its original
capitalized property names, its \code{Vector3D} return types and its
residue-building methods, kept intact so that proteusPy code runs
unchanged. \code{Vector3D} is the small vector class \code{Turtle3D}
uses, with helper functions for angles, dihedrals and distances. The
three modules total about 1{,}600 lines including docstrings, of which
\code{TurtleND} is about 500.

The tests cover every operation individually, the 200-step random
agreement check of Theorem~\ref{thm:compat}, the degenerate cases of
\op{Orient} and \op{OrientToward}, and the tape. Coverage is
held above 80 percent. The package requires Python 3.12 or 3.13,
is licensed BSD-3-Clause, and has a documentation site with a
frame-conventions page and an API reference at
\url{https://flux-frontiers.github.io/turtlend/}.

%======================================================================
\section{Conclusion}
\label{sec:conclusion}
%======================================================================

The turtle is a small object: a point, $N$ unit vectors, and one
rotation primitive. Its value is that every operation is expressed in
its own frame, so the description of a path, a molecule or a walk on a
manifold does not change when the ambient coordinates do. The
three-dimensional version has been placing atoms since 1990 and the
$N$-dimensional version is a strict generalization of it, with the
1990 conventions preserved so that the two agree to $10^{-7}$ after
200 random commands. The whole of the $N$-dimensional turtle fits in
about 500 lines of NumPy.

\bibliographystyle{plainnat}
\bibliography{turtlend}

\end{document}
